\section{总结与延伸}

\subsection{一条贯穿整场课的设计链}

本讲可以压缩成八步：先把任务写成动作条件未来预测；区分 translation 与 disambiguation；选择融合方向；构建空间对齐数据；用 masked modeling 学习条件分布；以 H\&E 等廉价模态做带校准的 imputation；用统一模型提出反事实；最后通过实验和失败数据更新世界模型。任何一步被省略，系统都可能从科学工具退化为漂亮演示。

\begin{importantbox}{核心结论}
多模态世界模型的价值不在“把更多数据塞进更大 Transformer”，而在于把不同传感器的互补信息组织成可检验的干预预测。模型必须知道何时能翻译、何时需要消歧、何时信息根本不可识别，并把不确定性带入实验选择。
\end{importantbox}

\subsection{工程上最值得带走的检查表}

\begin{enumerate}
  \item \textbf{任务：}预测目标是否对应真实决策，而非仅仅容易优化？
  \item \textbf{切分：}是否按患者、中心与时间切分，避免空间与批次泄漏？
  \item \textbf{融合：}新增模态究竟提供 translation 还是 disambiguation？
  \item \textbf{缺失：}训练时是否模拟部署时的整模态缺失与测量预算？
  \item \textbf{校准：}是否按基因、患者与域报告不确定性覆盖？
  \item \textbf{反事实：}是否用真实 perturbation 数据检验方向与效应大小？
  \item \textbf{闭环：}失败实验是否被记录并用于更新模型？
\end{enumerate}

\subsection{开放问题}

第一，如何在 patient--organ--tissue--cell--molecule 层级间动态分配计算，而不是固定压缩？第二，如何区分 H\&E 中可识别的分子状态与根本不可识别的状态？第三，如何让模型 counterfactual 与真实因果干预对齐？第四，如何在缺乏健康组织和随机治疗数据时建模选择偏差？第五，scientist agent 的价值应如何通过前瞻实验而不是文本 benchmark 评估？

这些问题共同指向一个更严格的 AI for Science 标准：模型不仅要生成合理答案，还要暴露证据来源、分布距离、预测区间和可失败的实验承诺。世界模型的终点不是一张更逼真的热图，而是一个能被现实持续纠正的研究系统。

\subsection{拓展阅读}

\begin{itemize}[itemsep=0.2em,topsep=0.2em]
  \item Stanford CS25 V5 课程页：\url{https://web.stanford.edu/class/cs25/past/cs25-v5/}
  \item 本讲官方录像：\url{https://www.youtube.com/watch?v=8kXIaUM3h1E}
  \item 方法关键词：CLIP（共享表示）、Masked Autoencoder（遮蔽学习）、cross-attention（定向融合）、adaptive LayerNorm（条件调制）。评估时优先寻找患者级留出、真实干预、校准和跨中心复现。
\end{itemize}

\end{document}
